% !TeX root = ../../Book3.tex
\section{极值长度}
\label{b3:sec:0803}

极值长度以所有非负密度为竞争者，比的是最短曲线长度的平方与面积。缩放密度不改变比值，因此它是一个没有欧氏单位的量。

\subsection{曲线族}
\label{b3:subsec:080301}

曲线族 \(\Gamma\) 是区域内局部可求长曲线的集合。本章圆环的径向族包含所有连接两端的曲线，允许以边界端作为极限端点；绕行族包含绕中心指数为1或 \(-1\) 的闭曲线。矩形族则连接指定两条相对边。

\ProseParagraphBreak
同一个区域配不同曲线族会给不同极值量。径向族和绕行族将得到互为倒数的数值。曲线参数化不重要，长度积分使用弧长；重复往返只增加非负长度，不会降低连接端点所需的最短量。允许曲线不光滑并不妨碍定义，只需按局部可求长性理解积分。

\subsection{可容许度量}
\label{b3:subsec:080302}

取 Borel 可测密度 \(\rho\ge0\)，要求
\[
 0<A_\rho=\int_U\rho^2\,\mathrm{d}A<\infty.
\]
定义 \(L_\rho(\gamma)=\int_\gamma\rho\,\mathrm{d}s\)，以及 \(L_\rho(\Gamma)=\inf_{\gamma\in\Gamma}L_\rho(\gamma)\)。零密度区域允许存在，个别曲线长度也允许无穷。极值问题只比较有限正面积的密度，不把 \(0/0\) 或 \(\infty/\infty\) 当作竞争值。

\ProseParagraphBreak
与“令每条曲线长度至少1，再最小化面积”的表述相比，本章不先使用可容许一词作为额外条件；若 \(L_\rho(\Gamma)>0\)，把密度除以这个数即可转成那种等价形式。若族里有任意短曲线使下确界为0，该密度的贡献就是0。

\ProseParagraphBreak
极值密度的整体倍数不影响商。将 \(\rho\) 乘 \(c>0\)，最短长度乘 \(c\)，面积乘 \(c^2\)，所以 \(L^2/A\) 不变。面积可以归一化为1，也可把最短长度归一化为1。在后一归一化中
\[
 \frac1{\lambda(\Gamma)}=\inf\left\{\int\rho^2\,\mathrm{d}A:
       \int_\gamma\rho\,|\mathrm{d}z|\ge1\quad(\gamma\in\Gamma)\right\},
\]
只在 \(0<\lambda<\infty\) 时直接写此互倒公式。它来自对每个密度作尺度调整，而不额外假设上确界达到。这个面积下确界也常被称为曲线族的模；不同文献的“模”可能指极值长度或其倒数，使用时必须先核对定义。

\ProseParagraphBreak
非负 Borel 密度可以在某些点为0或无穷，只要求总面积有限且非零。对零面积集合任意改变密度未必不影响每一条曲线的长度，例如一条恰落在该集合中的曲线；因此若把密度视为 \(L^2\) 等价类，就还要加入曲线族例外的处理。本章初等计算使用明确的密度代表，长度不等式按实际曲线验证，不把平面几乎处处相等自动转成所有路径的相等。

\subsection{极值长度}
\label{b3:subsec:080303}

\begin{definition}\label{b3:def:08-extremal-length}
\term[jizhi-changdu]{极值长度}[extremal length]定义为
\[
 \lambda_U(\Gamma)=\sup_\rho\frac{L_\rho(\Gamma)^2}{A_\rho}.
\]
\end{definition}
\begin{theorem}\label{b3:thm:08-extremal-conformal}
若 \(f:U\rightarrow V\) 双全纯，则 \(\lambda_U(\Gamma)=\lambda_V(f\Gamma)\)。若 \(\Gamma_1\subset\Gamma_2\)，则 \(\lambda(\Gamma_1)\ge\lambda(\Gamma_2)\)。
\end{theorem}
\begin{proof}
给目标密度 \(\sigma\)，拉回 \(\rho(z)=\sigma(f(z))|f'(z)|\)。链式法则给每条曲线长度相等，平面换元的 Jacobi 行列式为 \(|f'|^2\)，所以面积相等。逆映射把任意源密度也这样对应到目标密度，取上确界得相等。曲线族包含时，下确界在较小集合上较大，对每个密度比较后取上确界即得第二式。
\end{proof}
乘密度一个正数使长度与面积分别乘该数及其平方，故比值保持不变。真正重要的是密度的分布而非总尺度。

\begin{proposition}[圆环的 Dirichlet 能量]\label{b3:prop:08-annulus-capacity}
设 \(A=A(r,R)\)，\(u\) 为闭圆环上的 \(C^1\) 实函数，在内圆取0、外圆取1。则
\[
 \int_A|\nabla u|^2\,\mathrm{d}A\ge\frac{2\pi}{\log(R/r)}.
\]
等号由 \(u_0(z)=\log(|z|/r)/\log(R/r)\) 达到，并且在这类函数中唯一。
\end{proposition}
\begin{proof}
记 \(L=\log(R/r)\)。沿每条径向线，\(\int_r^R\partial_su\,\mathrm{d}s=1\)。Cauchy--Schwarz 给
\[
 1\le\left(\int_r^R |\partial_su|^2s\,\mathrm{d}s\right)
       \left(\int_r^R\frac{\mathrm{d}s}{s}\right).
\]
对角度积分并舍去非负角方向能量，得到下界 \(2\pi/L\)。\(u_0\) 的梯度为径向 \(1/(Ls)\)，直接积分给等号。等号必须使角导数为0，且各径向 Cauchy--Schwarz 等号使 \(\partial_su\) 为 \(c/s\)；边界值迫使 \(c=1/L\)，故只能为 \(u_0\)。
\end{proof}
这个最小能量称为圆环的容量，它恰是径向连接族极值长度的倒数，也恰是绕行族的极值长度。若把圆环切开并用 \(w=\log z\)，\(u_0\) 变成宽度为 \(L\) 的平直圆柱上的线性势。几何比较、曲线长度和能量三个观点在此给同一数值。

\ProseParagraphBreak
平面共形换元使实函数的 Dirichlet 能量不变：梯度模平方乘 \(|f'|^{-2}\)，面积元乘 \(|f'|^2\)，两项抵消。于是对光滑边界环域可通过环域映射把最小能量带过去；边界值的对应须先有合适的边界延拓。对于任意粗糙边界，能量边界条件需要另行定义，不能直接写“在每个边界点取0或1”并假定经典解存在。后面的势论将系统研究这种问题。

\subsection{环域计算}
\label{b3:subsec:080304}

\begin{theorem}\label{b3:thm:08-annulus-extremal}
在 \(A(r,R)\) 中，连接两端的族 \(\Gamma_{\mathrm r}\) 与绕行族 \(\Gamma_{\mathrm c}\) 满足
\[
 \lambda(\Gamma_{\mathrm r})=\frac{\log(R/r)}{2\pi},\qquad
 \lambda(\Gamma_{\mathrm c})=\frac{2\pi}{\log(R/r)}.
\]
\end{theorem}
\begin{proof}
取 \(\rho=1/|z|\)，面积为 \(2\pi\log(R/r)\)。任意径向连接曲线的长度至少 \(\log(R/r)\)，由沿曲线积分 \(|\mathrm{d}\log|z||\le|\mathrm{d}z|/|z|\) 得到；径向线达到此值，所以给第一式下界。任意绕行曲线长度至少 \(2\pi\)，由辐角总变化得到；圆周达到此值，给第二式下界。

任取密度 \(\sigma\)。径向线的长度 \(\ell(\theta)=\int_r^R\sigma(se^{\mathrm{i}\theta})\,\mathrm{d}s\) 由 Cauchy--Schwarz 满足
\[
 \ell(\theta)^2\le\log(R/r)\int_r^R\sigma(se^{\mathrm{i}\theta})^2s\,\mathrm{d}s.
\]
对角度平均，存在径向线使长度平方不超过 \(A_\sigma\log(R/r)/(2\pi)\)，得到第一式上界。对圆周长度 \(h(s)=\int_0^{2\pi}\sigma(se^{\mathrm{i}\theta})s\,\mathrm{d}\theta\)，有 \(h(s)^2\le2\pi s^2\int_0^{2\pi}\sigma^2\,\mathrm{d}\theta\)。再以 \(\mathrm{d}s/s\) 积分，存在半径使 \(h(s)^2\le2\pi A_\sigma/\log(R/r)\)，得到第二式上界。可测密度由 Tonelli 定理保证这些积分对几乎所有参数成立；取下确界不受零测参数集影响。
\end{proof}
极值密度的几何是把对数圆柱展开后的常密度。两族一纵一横，长度平方除面积互为倒数。

\begin{proposition}[试探势给极值长度下界]\label{b3:prop:08-test-potential}
设有界平面通道 \(Q\) 的两个指定边界段为 \(E_0,E_1\)，\(u\) 在其闭包附近 \(C^1\)，在两段分别取0、1，且 \(0<E=\int_Q|\nabla u|^2\,\mathrm{d}A<\infty\)。连接这两段的曲线族满足 \(\lambda(\Gamma)\ge1/E\)。
\end{proposition}
\begin{proof}
取密度 \(\rho=|\nabla u|\)。任一连接曲线由链式法则给 \(1=|u(\gamma(1))-u(\gamma(0))|\le\int_\gamma\rho\,|\mathrm{d}z|\)，所以最短长度至少1，总面积为 \(E\)，代入定义即得下界。
\end{proof}
圆环的对数势达到最佳下界，矩形的线性势也达到最佳下界。一般通道可先构造一个方便的 \(u\)，即使它不是调和函数，仍给严格的下界。越小的能量给越强的下界；寻找最小能量自然引向 Dirichlet 问题，但这个简单比较本身不需要先证明最小解存在。若势只具有弱导数，沿所有曲线的链式法则需再处理，不能把本命题的 \(C^1\) 条件无声地删除。
